\subsection{投射模}

定义1. 一个A-模P称为投射的，如果对每个由A-线性映射组成的正合序列 \(\mathbf{F}' \to \mathbf{F} \to \mathbf{F}''\) ，序列

\[
\operatorname{Hom} (\mathrm{P}, \mathrm{F} ^ {\prime}) \rightarrow \operatorname{Hom} (\mathrm{P}, \mathrm{F}) \rightarrow \operatorname{Hom} (\mathrm{P}, \mathrm{F} ^ {\prime \prime})
\]

是正合的。

命题3. 对于A-模P，一族子模 \((\mathbf{M}_{t})\) 的直和是投射的，当且仅当每一个 \(\mathbf{M}_{t}\) 都是投射的。

对于每个A-模同态 \(u: \mathbf{E} \to \mathbf{F}\) ,

\[
\operatorname{Hom} \left(1 _ {\mathrm{P}}, u\right): \operatorname{Hom} (\mathrm{P}, \mathrm{E}) \rightarrow \operatorname{Hom} (\mathrm{P}, \mathrm{F})
\]

与 \(\prod_{i}\operatorname{Hom}(1_{M_i},u)\) 等同（§1，第6号，命题6的推论1）；因此，该结论由定义1及§1，第5号，命题5（ii）得出。

推论。每个自由A-模都是投射的。

由命题3，只需证明 \(A_{s}\) 是投射的，而这由§1第14号中图表（50）的交换性立即得出。

命题4. 设P是一个A-模。以下性质是等价的：

\begin{enumerate}
\def\labelenumi{(\alph{enumi})}
\item
  P 是投射的。
\item
  对于每一个由A-线性映射组成的正合序列 \(0 \to \mathbf{F}' \to \mathbf{F} \to \mathbf{F}'' \to 0\) ，序列
\end{enumerate}

\[
0 \rightarrow \operatorname{Hom} (\mathrm{P}, \mathrm{F} ^ {\prime}) \rightarrow \operatorname{Hom} (\mathrm{P}, \mathrm{F}) \rightarrow \operatorname{Hom} (\mathrm{P}, \mathrm{F} ^ {\prime \prime}) \rightarrow 0
\]

是正合的。

\begin{enumerate}
\def\labelenumi{(\alph{enumi})}
\setcounter{enumi}{2}
\item
  对于每个满射的A-模同态 \(u: \mathbf{E} \to \mathbf{E}''\) 以及每个同态 \(f: \mathbf{P} \to \mathbf{E}''\) ，存在一个同态 \(g: \mathbf{P} \to \mathbf{E}\) 使得 \(f = u \circ g\) （即称 \(f\) 可以“提升”为 \(\mathbf{P}\) 到 \(\mathbf{E}\) 的一个同态）.
\item
  每个由A-线性映射组成的正合序列 \(0 \to \mathbf{E}' \to \mathbf{E} \xrightarrow{v} \mathbf{P} \to 0\) 都分裂（因此 \(\mathbf{P}\) 同构于 \(\mathbf{E}\) 的一个直因子）.
\item
  P同构于自由A-模的一个直因子。
\end{enumerate}

(a) 蕴含 (b) 是显然的。要看出 (b) 蕴含 (c)，只需将 (b) 应用于正合序列 \(0 \to \mathbf{E}' \to \mathbf{E} \xrightarrow{u} \mathbf{E}'' \to 0\) ，其中 \(\mathbf{E}' = \operatorname{Ker}(u)\) ，因为 (c) 表达了如下事实

\[
\operatorname{Hom} \left(\mathbf {1} _ {\mathrm{P}}, u\right): \operatorname{Hom} (\mathrm{P}, \mathrm{E}) \rightarrow \operatorname{Hom} (\mathrm{P}, \mathrm{E} ^ {\prime \prime})
\]

是满射的。为了证明 (c) 蕴含 (d)，只需将 (c) 应用于满射同态 \(v: \mathbf{E} \to \mathbf{P}\) 及同态 \(\mathbf{l}_{\mathbf{P}}: \mathbf{P} \to \mathbf{P}\) ；存在一个同态 \(g: \mathbf{P} \to \mathbf{E}\) 使得 \(\mathbf{l}_{\mathbf{P}} = v \circ g\) ，这一事实意味着序列

\[
0 \longrightarrow \mathrm{E} ^ {\prime} \longrightarrow \mathrm{E} \stackrel {v} {\longrightarrow} \mathrm{P} \longrightarrow 0
\]

分裂（§1，第9号，命题15）。由于对每个A-模M，存在一个自由A-模L及一个正合序列 \(0 \to R \to L \to M \to 0\) (§ 1，第11号，命题20)，显然(d)蕴含(e)。最后，由命题3及其推论可知，(e)蕴含(a)。

推论1. 一个A-模是投射的且有限生成的，当且仅当它是具有有限基的自由A-模的一个直因子。

该条件显然是充分的；反之，一个有限生成的投射模E同构于一个具有有限基的自由模F的商模（§1，第11号），且根据命题4(d)，E同构于F的一个直因子。

推论2. 设C为交换环，E、F为两个有限生成的投射C-模；则 \(\operatorname{Hom}_{\mathbf{C}}(\mathbf{E}, \mathbf{F})\) 是一个有限生成的投射C-模。

可以假设存在两个有限生成的自由C-模M，N，使得 \(M = E \oplus E'\) , \(N = F \oplus F'\) ；由§1，第6号，命题6的推论1可知 \(\operatorname{Hom}_{\mathbb{C}}(\mathbf{M}, \mathbf{N})\) 是有限生成的且自由的，而另一方面 \(\operatorname{Hom}_{\mathbb{C}}(\mathbf{M}, \mathbf{N})\) 同构于

\[
\operatorname{Hom} _ {\mathbf {c}} (\mathbf {E}, \mathbf {F}) \oplus \operatorname{Hom} _ {\mathbf {c}} (\mathbf {E} ^ {\prime}, \mathbf {F}) \oplus \operatorname{Hom} _ {\mathbf {c}} (\mathbf {E}, \mathbf {F} ^ {\prime}) \oplus \operatorname{Hom} (\mathbf {E} ^ {\prime}, \mathbf {F} ^ {\prime}),
\]

由此得出推论。

